%%correc_saut_mnq
A chaque instant $j$ :
\begin{itemize}
\item un saut est d\'etect\'e si : \newline
$X[j] - X[j-1] > S_0$ et $X[j] - X[j+1] > S_0$  (pic vers le haut)

ou 

$X[j-1] - X[j] > S_0$ et $X[j+1] - X[j] > S_0$  (pic vers le bas)

\hspace*{1cm} alors $flag[j] = 1$ et $X_c[j] = \frac{X[j+1] + X[j-1]}{2}$ \newline

\item un saut apr�s une valeur manquante est d\'etect\'e si :\newline
$X[j-1] = mnq$
et
$|X[j+1] - X[j]| > S_{mnq}$ 

\hspace*{1cm} alors $flag[j] = 2$ et $X_c[j] = mnq$ \newline

\item un saut avant une valeur manquante est d\'etect\'e si :\newline
$X[j+1] = mnq$
et
$|X[j] - X[j-1]| > S_{mnq}$ 

\hspace*{1cm} alors $flag[j] = -2$ et $X_c[j] = X[j]$ \newline

\item une forte variation positive est d\'etect\'e si :\newline
$X[j+1] - X[j] > S_0$ et $X[j] - X[j-1] > S_0$  (forte variation vers le haut)

\hspace*{1cm} alors $flag[j] = -1$ et  $X_c[j] = X[j]$ \newline

\item une forte variation n�gative est d\'etect\'e si :\newline
$X[j-1] - X[j] > S_0$ et $X[j] - X[j+1] > S_0$  (forte variation vers le bas)

\hspace*{1cm} alors $flag[j] = -1$ et $X_c[j] = X[j]$ \newline

\item aucune variation "particuli�re" n'est d�tect�e :\newline
\hspace*{1cm} alors $flag[j] = 0$ et $X_c[j] = X[j]$ \newline

\end{itemize}
